\documentclass[12pt, a4paper]{article}
\usepackage[utf8]{inputenc}
\usepackage{geometry}
\geometry{a4paper, margin=1in}
\usepackage{graphicx}

\usepackage{setspace}
\usepackage{titlesec}
\usepackage{booktabs}
\usepackage{amsmath}
\usepackage{amssymb}
\usepackage{enumitem}
\usepackage{tcolorbox} % For highlighting progress

\setstretch{1.2}

\titleformat{\section}{\Large\bfseries}{\thesection}{1em}{}
\titleformat{\subsection}{\large\bfseries}{\thesubsection}{1em}{}
\titleformat{\subsubsection}{\normalsize\itshape}{\thesubsubsection}{1em}{}

\usepackage[colorlinks=true, allcolors=blue]{hyperref}

\begin{document}

\begin{titlepage}
    \centering
    \vspace*{2cm}
    
    {\Huge \textbf{Team Handbook \& Internal Guide}} \\[1cm]
    {\LARGE \textbf{Architecture-Agnostic Layer-Wise PTQ of LLMs}} \\[0.5cm]
    {\large \textbf{B.Tech Major Project 2026--27}} \\[2cm]
    
    \textbf{Team Members:} \\[0.2cm]
    \begin{tabular}{ll}
        Faqre Alam & (23/CS/150) \\
        Ekansh Agrawal & (23/CS/149) \\
        Guneet Toppo & (23/CS/160)
    \end{tabular} \\[2cm]
    
    \begin{tcolorbox}[colback=blue!5!white,colframe=blue!75!black,title=Project Status Tracker]
        \textbf{Current Status (As of Sep 2026):} We have successfully built the framework and are executing \textbf{Phase 1} (The 27-run ablation grid on Kaggle). The mathematical groundwork (\textbf{Phase 0}) is complete. \textbf{Phase 2} (Novel Algorithm) and \textbf{Phase 3} (Writing) are up next.
    \end{tcolorbox}
    
    \vfill
    {\large Delhi Technological University}
\end{titlepage}

\tableofcontents
\newpage

\section{Introduction: The Memory Wall}
The core problem our project solves is the \textbf{Memory Wall}. When running LLMs (like generating text token-by-token), the bottleneck is not computation (FLOPs) but memory bandwidth. For every single token generated, the entire weight matrix of the LLM must be loaded from memory to the GPU cores. 

By applying \textbf{Post-Training Quantization (PTQ)}, we compress the model weights from 16-bit down to 4-bit or 2-bit. This directly shrinks the memory footprint by $4\times$ to $8\times$, drastically speeding up the model and allowing large models to run on small edge devices.

\section{The Three Core PTQ Methods We Are Testing}
To understand what we are building, we must understand the three algorithms implemented in our framework:

\subsection{1. Round-To-Nearest (RTN) --- Our Baseline}
This is the simplest method. We simply find the maximum and minimum values of a weight channel and map all 16-bit float values to the nearest integer grid (e.g., 16 levels for 4-bit). 
\begin{itemize}
    \item \textbf{Pros:} Very fast, requires no calibration data.
    \item \textbf{Cons:} Disastrous at 2-bit. It completely destroys the LLM because it introduces massive rounding error to sensitive weights.
\end{itemize}

\subsection{2. GPTQ --- Second-Order Error Compensation}
GPTQ uses \textbf{calibration data} (128 sequences from WikiText-2). When we quantize a weight column and introduce an error, GPTQ uses a mathematical matrix called the \textbf{Hessian} ($H = 2 X^T X$) to calculate exactly how to adjust the remaining unquantized columns to cancel out that error.
\begin{itemize}
    \item \textbf{Result:} It provides near-lossless compression at 4-bit, but starts struggling at 2-bit due to ``outlier'' weights.
\end{itemize}

\subsection{3. QuIP/Hadamard --- The 2-bit Enabler}
Before applying GPTQ, we multiply the weight matrix by a randomized \textbf{Hadamard Rotation Matrix}. This mathematical trick smears outlier weights evenly across all columns. 
\begin{itemize}
    \item \textbf{Result:} It equalizes the weight distribution, allowing GPTQ to work much better at 2-bit (about 13\% error reduction over standard GPTQ).
\end{itemize}

\section{The 27-Run Architecture Ablation Study (Phase 1)}
We are testing if the above methods work equally well on different types of LLM architectures. Our framework dynamically injects \texttt{nn.Linear} hooks to support any architecture without custom code.

We are running a massive 27-run evaluation grid on Kaggle T4 GPUs:
\begin{itemize}
    \item \textbf{3 Models:} Qwen2.5-0.5B (GQA 7:1), Llama-3.2-1B (GQA 4:1), Pythia-1.4B (Full MHA)
    \item \textbf{3 Methods:} RTN, GPTQ, GPTQ+Hadamard
    \item \textbf{3 Bit-widths:} 2-bit, 3-bit, 4-bit
\end{itemize}
\textbf{Status:} \textit{Currently executing. This will produce the core ablation table for our final report.}

\section{Our Novel Contribution: Hessian-Diagonal Allocation (Phase 2)}
Once Phase 1 is done, we will build our own algorithm. Instead of quantizing every column to exactly 3 bits, we will allocate \textbf{4 bits to sensitive columns} and \textbf{2 bits to insensitive columns}, ensuring the average is exactly 3 bits.

We identify sensitivity using our novel, cheap formulation derived from the Taylor expansion of the loss function:
\begin{equation}
    \hat{s}_j = \mathbf{H}_{jj} \cdot \mathrm{Var}(\mathbf{W}_{:,j})
\end{equation}
Where $\mathbf{H}_{jj}$ is the diagonal of the Hessian (captured during calibration), and $\mathrm{Var}$ is the variance of the weight column. We will use a reverse-water-filling approach to allocate bits based on this score.

\section{Roadmap and Timeline}

\begin{tcolorbox}[colback=green!5!white,colframe=green!60!black,title=Phase 0 \& 1: The Foundation (Completed \& Active)]
\textbf{Goal:} Build the universal quantization pipeline and run baselines.\\
\textbf{What we did:} 
\begin{itemize}[leftmargin=*]
    \item We initialized the PyTorch pipeline.
    \item We fixed memory issues with Apple Silicon (MPS).
    \item We created the \texttt{run.py} entry point and hooked up the \texttt{lm-evaluation-harness}.
    \item \textbf{Current Action:} Waiting for the Kaggle runs to finish producing the \texttt{ablation\_table.csv}.
\end{itemize}
\end{tcolorbox}

\begin{tcolorbox}[colback=yellow!5!white,colframe=yellow!70!black,title=Phase 2: The Novelty (Upcoming Jan--Feb 2027)]
\textbf{Goal:} Implement the Hessian-Diagonal Mixed Precision.\\
\textbf{Tasks:} 
\begin{itemize}[leftmargin=*]
    \item Code the $\hat{s}_j$ sensitivity algorithm in PyTorch.
    \item Run a sweep varying the ratio of 4-bit vs 2-bit weights.
    \item Prove that our method achieves a lower Perplexity (PPL) than the uniform 3-bit baseline from Phase 1.
\end{itemize}
\end{tcolorbox}

\begin{tcolorbox}[colback=red!5!white,colframe=red!70!black,title=Phase 3: Synthesis \& Publication (Upcoming Mar--May 2027)]
\textbf{Goal:} Write the final BTech report and potentially publish it.\\
\textbf{Tasks:} 
\begin{itemize}[leftmargin=*]
    \item Generate graphs plotting Perplexity vs. Average Bits.
    \item Write the paper targeting ACL SRW or arXiv.
    \item Prepare final presentation slides for the viva voce.
\end{itemize}
\end{tcolorbox}

\end{document}
